\documentclass[border=6pt]{standalone}
\usepackage[utf8]{inputenc}
\usepackage[T5]{fontenc}
\usepackage{tikz}
\usetikzlibrary{shapes.multipart, arrows.meta}

\begin{document}
\begin{tikzpicture}[
    font=\small,
    cls/.style={rectangle split, rectangle split parts=3, draw=blue!60!black, thick,
                rectangle split part fill={blue!15, white, white},
                text width=#1, align=left, inner sep=4pt},
    cls/.default=5.6cm,
    simple/.style={rectangle, draw=blue!60!black, thick, fill=blue!8,
                   minimum width=1.55cm, minimum height=0.7cm},
    inh/.style={-{Triangle[open, length=3mm, width=3mm]}, thick}
]
    \node[cls, anchor=north] (base) at (0, 0) {
        \centering \textbf{\textit{Block}}
        \nodepart{second}
        + shape : char[4][4]
        \nodepart{third}
        + \textit{rotate()} : void\\
        + \textit{undoRotate()} : void\\
        + \textit{$\sim$Block()}
    };

    \node[cls=4.3cm, anchor=north] (I) at (-7.0, -3.6) {
        \centering \textbf{IBlock}
        \nodepart{second}
        - isVertical : bool
        \nodepart{third}
        + rotate() : void\\
        + undoRotate() : void
    };
    \node[cls=4.3cm, anchor=north] (O) at (7.0, -3.6) {
        \centering \textbf{OBlock}
        \nodepart{second}
        \phantom{x}
        \nodepart{third}
        + rotate() : void (rỗng)\\
        + undoRotate() : void (rỗng)
    };

    \node[simple, anchor=north] (J) at (-3.5, -3.6) {JBlock};
    \node[simple, anchor=north] (L) at (-1.75, -3.6) {LBlock};
    \node[simple, anchor=north] (S) at ( 0.0, -3.6) {SBlock};
    \node[simple, anchor=north] (T) at ( 1.75, -3.6) {TBlock};
    \node[simple, anchor=north] (Z) at ( 3.5, -3.6) {ZBlock};

    \coordinate (bus) at (0, -2.9);
    \draw[thick] (I.north) |- (bus) -| (O.north);
    \foreach \n in {J, L, T, Z} \draw[thick] (\n.north) -- (\n.north |- bus);
    \draw[thick] (S.north) -- (bus);
    \draw[inh] (bus) -- (base.south);

    \node[font=\footnotesize, text=black!70, align=center] at (0, -4.75)
        {J, L, S, T, Z dùng cách xoay mặc định của \textit{Block}};
    \node[font=\footnotesize, text=black!70, anchor=north west] at (-9.2, 0)
        {\textit{In nghiêng}: lớp trừu tượng / hàm ảo};
\end{tikzpicture}
\end{document}
